\section{Conics: complete parametrization and explicit height growth}
\label{sec:conics}

The first substantial geometric application is a complete description of
the rational points on a nonsingular conic \emph{with a supplied rational
point}.  The geometric idea is projection from that point: a line through
it has one other intersection with the conic, with the tangent line
returning the original point.  This is the usual presentation of rational
points on a conic; see, for example, \cite{sutherland2025}.  We give all the
formulas, handle the tangent case explicitly, and derive quantitative
height estimates by integer identities.  Consequently neither a theorem
about algebraically closed fields nor an intersection-multiplicity
formalism is needed.

\subsection{A normal form that preserves explicit information}

Let
\[
 q(X_0,X_1,X_2)=\sum_{0\leq i\leq j\leq2}a_{ij}X_iX_j,
 \qquad a_{ij}\in\Z,
\]
and let
\[
 S_q=\begin{pmatrix}
 2a_{00}&a_{01}&a_{02}\\
 a_{01}&2a_{11}&a_{12}\\
 a_{02}&a_{12}&2a_{22}
 \end{pmatrix}.
\]
We use the decidable nonsingularity condition $\det S_q\ne0$.
It implies the usual assertion that the projective conic has no singular
point: its gradient at a vector $x$ is $S_qx$, and multiplication by the
adjugate matrix shows that $S_qx=0$ implies $x=0$.  All matrices here have
the fixed size three, so this argument is an identity of integer
polynomials.  Put
\[
 C=\sum_{i\leq j}|a_{ij}|.
\]
Nonsingularity implies $C\geq1$.

Suppose that a nonzero primitive integer vector
$p=(p_0,p_1,p_2)$ with $q(p)=0$ is given, and write
$P=\max_i|p_i|\geq1$.  A permutation of coordinates makes $p_0\ne0$.
This permutation is found by checking the three coordinates; it preserves
$C$ and $P$.  Define
\[
 M=\begin{pmatrix}p_0&0&0\\p_1&1&0\\p_2&0&1\end{pmatrix}.
\]
Since $\det M=p_0\ne0$, it induces a bijection of rational projective
planes.  Its inverse can be evaluated without rational division:
\begin{equation}\label{eq:conic-adjugate}
 \operatorname{adj}(M)x
 =(x_0,\ p_0x_1-p_1x_0,\ p_0x_2-p_2x_0).
\end{equation}
Indeed $M\operatorname{adj}(M)=p_0I$.

The coefficient of $U^2$ in $q(M(U,V,W))$ is $q(p)=0$.
Write the remaining coefficients as
\begin{equation}\label{eq:conic-normal}
 q(M(U,V,W))=U L(V,W)+R(V,W),
 \qquad
 L(V,W)=\ell V+mW,\quad R(V,W)=aV^2+bVW+cW^2.
\end{equation}
Explicitly,
\[
\begin{split}
 \ell&=a_{01}p_0+2a_{11}p_1+a_{12}p_2,\\
 m&=a_{02}p_0+a_{12}p_1+2a_{22}p_2,\\
 a&=a_{11},\qquad b=a_{12},\qquad c=a_{22}.
\end{split}
\]
The polar matrix of the transformed form is $M^{\mathsf T}S_qM$.
Expansion of its determinant gives
\begin{equation}\label{eq:conic-r}
 \det(M^{\mathsf T}S_qM)
 =\det\begin{pmatrix}0&\ell&m\\\ell&2a&b\\m&b&2c\end{pmatrix}
 =-2r,
 \qquad r=am^2-b\ell m+c\ell^2=R(-m,\ell).
\end{equation}
Since this determinant also equals $p_0^2\det S_q$, we have $r\ne0$.
In particular $(\ell,m)\ne(0,0)$.

The following two identities will serve both to exclude undefined cases
and to control the cancellation in projective coordinates:
\begin{align}
 rs^2&=m^2R(s,t)
   -L(s,t)\bigl((bm-c\ell)s+cmt\bigr),\label{eq:conic-identity-s}\\
 rt^2&=\ell^2R(s,t)
   -L(s,t)\bigl(a\ell s+(b\ell-am)t\bigr).
   \label{eq:conic-identity-t}
\end{align}
For example, expansion of the right side of
\eqref{eq:conic-identity-s} cancels its $st$ and $t^2$ terms and leaves
$(am^2-b\ell m+c\ell^2)s^2$.  Expanding the other identity gives the
same verification.  If $L(s,t)=R(s,t)=0$, these identities and $r\ne0$
imply $s=t=0$.  Thus $L$ and $R$ have no common projective zero over
$\Q$; this conclusion has been obtained without a resultant theorem.

\subsection{The parametrization, including its tangent value}

Write $C_q(\Q)$ for the rational projective points satisfying $q=0$.
For $(s,t)\ne(0,0)$ define
\begin{equation}\label{eq:conic-param}
 F(s,t)=\bigl(-R(s,t),\ sL(s,t),\ tL(s,t)\bigr),
 \qquad
 \Phi([s:t])=[M F(s,t)].
\end{equation}
All projective vectors are normalized by the primitive normalization
procedure when a canonical point code is required.

\begin{theorem}[Effective parametrization of a pointed conic]
\label{thm:conic-param}
Under the hypotheses above, $\Phi$ is a bijection
\[
 \Phi:\PP^1(\Q)\longrightarrow C_q(\Q).
\]
It and its inverse are primitive recursive on canonical point codes,
uniformly in the six coefficients and the supplied point $p$.
The parameter $[-m:\ell]$ maps to $[p]$.
\end{theorem}

\begin{proof}
First, $F(s,t)\ne0$.  If $L(s,t)\ne0$, at least one of its last two
coordinates is nonzero.  If $L(s,t)=0$, then $R(s,t)\ne0$ by
\eqref{eq:conic-identity-s}--\eqref{eq:conic-identity-t}.
Since $M$ is invertible, $MF(s,t)\ne0$ as well.  All coordinates of $F$
are homogeneous of degree two.  Rescaling $(s,t)$ by a nonzero rational
number rescales $F$ by its square; hence the formula is well defined
on projective points.

Substitution in \eqref{eq:conic-normal}, using homogeneity of $L$ and
$R$, gives
\[
\begin{aligned}
 q(MF(s,t))&=(-R(s,t))L(sL(s,t),tL(s,t))\\
           &\qquad+R(sL(s,t),tL(s,t))\\
           &=-R(s,t)L(s,t)^2+R(s,t)L(s,t)^2=0.
\end{aligned}
\]
Thus the image lies on the conic.

To construct the inverse, first transform a conic point $[x]$ by
\eqref{eq:conic-adjugate}, obtaining $[u:v:w]$ satisfying
\[
 uL(v,w)+R(v,w)=0.
\]
If $(v,w)\ne(0,0)$, then $L(v,w)\ne0$: otherwise this equation would
also give $R(v,w)=0$, contrary to the two identities above.  Therefore
\[
 F(v,w)=(-R(v,w),vL(v,w),wL(v,w))
       =L(v,w)(u,v,w).
\]
In this case the inverse parameter is exactly $[v:w]$.
If $(v,w)=(0,0)$, the point is $[u:0:0]=[1:0:0]$ in the transformed
coordinates, and hence $[x]=[p]$.  Its inverse parameter is defined to
be $[-m:\ell]$.  This pair is nonzero, its $L$-value is zero, and its
$R$-value is $r\ne0$, so it indeed maps to $[1:0:0]$.

It remains to check that the inverse just specified recovers every
parameter.  When $L(s,t)\ne0$, the last two coordinates of $F(s,t)$
recover $[sL:tL]=[s:t]$.  When $L(s,t)=0$, the nonzero vector $(s,t)$
is proportional to $(-m,\ell)$.  For completeness, if $\ell\ne0$,
the equation $\ell s=-mt$ gives
$(s,t)=(t/\ell)(-m,\ell)$ and forces $t\ne0$; if $\ell=0$, then
$m\ne0$, $t=0$, and both projective pairs equal $[1:0]$.
Thus this exceptional case also recovers the original parameter.

Every branch tests an integer equality or inequality.  Matrix products,
polynomial values, signed integer arithmetic, gcd, and exact division
are primitive recursive.  The displayed inverse uses integer adjugate
coordinates followed by normalization, so rational division in the
explanatory proportionality check is not an operation required by the
algorithm.  The formulas define total programs by returning a fixed
default value on inputs failing the decidable hypotheses.  The theorem
asserts their correctness on valid inputs.
\end{proof}

The exceptional value has a geometric interpretation that is already
visible in the arithmetic.  At $(U,V,W)=(1,0,0)$ the gradient of the
transformed conic is $(0,\ell,m)$, so its tangent is
$\ell V+mW=0$.  This is exactly the line represented by $[-m:\ell]$.
Thus the tangent is included in the bijection rather than discarded as
an undefined slope.

\subsection{Quadratic height comparison with computable constants}

For a rational projective point, $\Ht$ denotes the maximum absolute
coordinate of its primitive integer representative.  Define the
following positive integers from the coefficients and the seed:
\begin{equation}\label{eq:conic-height-constants}
\begin{split}
 D_s&=m^2+|bm-c\ell|+|cm|,\\
 D_t&=\ell^2+|a\ell|+|b\ell-am|,\\
 D&=\max(D_s,D_t),\qquad A=3CP,\qquad K=2PD.
\end{split}
\end{equation}
Here $D\geq1$ because $(\ell,m)\ne(0,0)$.

\begin{theorem}[Height of a conic parametrization]
\label{thm:conic-height}
For every $T\in\PP^1(\Q)$,
\begin{equation}\label{eq:conic-height}
 \Ht(T)^2\leq K\Ht(\Phi(T)),
 \qquad
 \Ht(\Phi(T))\leq A\Ht(T)^2.
\end{equation}
In particular, a complete list of conic points of height at most $B$
is obtained by evaluating $\Phi$ at all canonical parameters of height
at most $\lfloor\sqrt{KB}\rfloor$ and retaining exactly those images
of height at most $B$.
\end{theorem}

\begin{proof}
Take a primitive integer pair $(s,t)$ representing $T$, and put
$S=\max(|s|,|t|)\geq1$.
The coefficient formulas give
\[
 |R(s,t)|\leq CS^2,
 \qquad |\ell|+|m|\leq2CP,
 \qquad |L(s,t)|\leq2CPS.
\]
The coordinates of $MF$ are
\[
 -p_0R,\qquad -p_1R+sL,\qquad -p_2R+tL.
\]
Their absolute values are at most $3CPS^2=AS^2$.  Primitive
normalization cannot increase that maximum.  This proves the upper
inequality.

For the lower inequality it is essential to control cancellation of
common factors.  Let
\[
 g=\gcd\bigl(|R(s,t)|,|sL(s,t)|,|tL(s,t)|\bigr).
\]
This integer is positive because $F\ne0$.
Since $\gcd(s,t)=1$, a B\'ezout identity for $s,t$ shows
$\gcd(sL,tL)=|L|$, including the case $L=0$.  Consequently
$g=\gcd(|R|,|L|)$.  The identities
\eqref{eq:conic-identity-s}--\eqref{eq:conic-identity-t} show that
$g$ divides $rs^2$ and $rt^2$.  Since $s^2,t^2$ are coprime, another
B\'ezout identity yields $g\mid r$.  Hence $g\leq|r|$.

Put $T_0=\max(|R|,S|L|)=\max_i|F_i|$.
If $|s|=S$, \eqref{eq:conic-identity-s} gives
\[
 |r|S^2\leq m^2|R|+
 (|bm-c\ell|+|cm|)S|L|\leq D_sT_0\leq DT_0.
\]
If $|t|=S$, the same conclusion follows from
\eqref{eq:conic-identity-t} and $D_t$.
Since $T_0=g\Ht([F])$ and $g\leq|r|$, cancellation of the positive
integer $|r|$ gives
\begin{equation}\label{eq:conic-intermediate-height}
 S^2\leq D\Ht([F]).
\end{equation}
Now let $x$ be the primitive representative of $\Phi(T)=[MF]$.
The integer vector $\operatorname{adj}(M)x$ is nonzero and represents
$[F]$.  Formula \eqref{eq:conic-adjugate} bounds each of its
coordinates by $2P\Ht(\Phi(T))$.  Normalizing it and substituting in
\eqref{eq:conic-intermediate-height} yields
$S^2\leq2PD\Ht(\Phi(T))=K\Ht(\Phi(T))$.

For the last assertion, any image of height at most $B$ has inverse
parameter with squared height at most $KB$.  The parameter box is
therefore complete, and its final height filter removes exactly the
unwanted images.  The integer square root is obtained by a bounded
search through $0,\ldots,KB$.
\end{proof}

The theorem also supplies an inverse height estimate in its strongest
useful elementary form:
\[
 \Ht(\Phi^{-1}(X))\leq\lfloor\sqrt{K\Ht(X)}\rfloor.
\]
The linear bound obtained directly from the adjugate is easier, but
would lose the correct quadratic relation between the parameter
height and the conic height.

\subsection{A finite counting theorem and constructive infinitude}

Let $N_q(B)$ be the number of canonical points of $C_q(\Q)$ with
height at most $B$.  This is a primitive recursive finite count.  In
the next theorem, inequalities with rational-looking constants are
stated after clearing denominators.

\begin{lemma}[Many coprime pairs in a square]
\label{lem:conic-coprime-count}
Let $G(N)$ count pairs $(s,t)$ with $1\leq s,t\leq N$ and
$\gcd(s,t)=1$.  For every $N\geq1$,
\[
 4G(N)\geq N^2.
\]
\end{lemma}

\begin{proof}
For $N=1$ the assertion is immediate.  For $N\geq2$, every pair
with gcd greater than one belongs to at least one of the sets
\[
 E_d=\{(s,t):1\leq s,t\leq N,\ d\mid s,\ d\mid t\},
 \qquad 2\leq d\leq N.
\]
The finite union bound and $\#E_d=\lfloor N/d\rfloor^2$ give
\[
 N^2-G(N)\leq\sum_{d=2}^N\lfloor N/d\rfloor^2
 \leq N^2\sum_{d=2}^N\frac1{d^2}.
\]
There is no infinite series in this argument.  For $d\geq3$,
$1/d^2\leq1/(d(d-1))$, and the finite telescoping identity gives
\[
 \sum_{d=2}^N\frac1{d^2}
 \leq\frac14+\sum_{d=3}^N\left(\frac1{d-1}-\frac1d\right)
 =\frac34-\frac1N.
\]
Thus $G(N)\geq N^2/4+N$, which in particular proves the stated
inequality.  Every displayed rational comparison can equivalently
be multiplied by the positive product of its denominators.
\end{proof}

\begin{theorem}[Explicit linear growth of the point count]
\label{thm:conic-count}
With $A,K$ from \eqref{eq:conic-height-constants},
\[
 N_q(B)\leq4KB\quad(B\geq1),
 \qquad
 B\leq16A N_q(B)\quad(B\geq A).
\]
For each $n\geq1$ one can also produce $n+1$ distinct points on
the conic, each of height at most $An^2$.
\end{theorem}

\begin{proof}
For the upper bound put $S=\lfloor\sqrt{KB}\rfloor\geq1$.
The inverse of every point counted by $N_q(B)$ is a parameter of
height at most $S$.  The number of all nonzero integer pairs in the
box $[-S,S]^2$ is $4S^2+4S$.  Exactly one element in each pair
$\{v,-v\}$ has its first nonzero coordinate positive.  Thus there
are at most $2S^2+2S\leq4S^2$ canonical primitive pairs in this box.
The bijection and $S^2\leq KB$ imply $N_q(B)\leq4KB$.

For the lower bound choose the largest integer $N$ with $AN^2\leq B$.
The condition $B\geq A$ ensures $N\geq1$.
Every positive coprime pair in $[1,N]^2$ gives a different
projective parameter: two primitive positive integer pairs that
are rational multiples of one another are equal.  Their images
are distinct by Theorem~\ref{thm:conic-param} and have height at
most $AN^2\leq B$.  Hence
\[
 4N_q(B)\geq4G(N)\geq N^2.
\]
Maximality of $N$ gives
$B<A(N+1)^2\leq4AN^2\leq16AN_q(B)$, proving the assertion.

For the last statement use the pairwise distinct parameters
$[k:1]$ for $k=0,1,\ldots,n$.  Their heights are at most $n$ when
$n\geq1$, and Theorem~\ref{thm:conic-height} supplies the stated
bound.  This yields an explicit list for each $n$ and therefore
the constructive content of infinitude.
\end{proof}

This is a height-counting result rather than an appeal to an
asymptotic counting theorem.  Both constants and the enumerating
algorithm are given by fixed primitive recursive expressions.
The lower estimate already captures the linear order of growth
of the number of rational points on a pointed conic with respect
to its projective height.

\subsection{Two fully arithmetic examples}

\paragraph{The unit circle.}
For $X^2+Y^2=Z^2$, the preceding construction has the familiar form
\begin{equation}\label{eq:circle-param}
 [s:t]\longmapsto[s^2-t^2:2st:s^2+t^2].
\end{equation}
Its defining identity is
$(s^2-t^2)^2+(2st)^2=(s^2+t^2)^2$.
An inverse is $[Z+X:Y]$ unless $(X:Y:Z)=(-1:0:1)$; that point has
inverse $[0:1]$.  Indeed if $Z+X=0$, the conic equation gives
$Y^2=0$ and hence $Y=0$, so the exceptional point is the only
undefined value of the displayed inverse pair.  Otherwise
substitution into \eqref{eq:circle-param}, followed by use of
$X^2+Y^2=Z^2$, recovers $[X:Y:Z]$.

If $\gcd(s,t)=1$, the common divisor of the three coordinates in
\eqref{eq:circle-param} is exactly
\[
 \epsilon=\begin{cases}2,&s,t\text{ both odd},\\1,&s,t\text{ of opposite parity}.
 \end{cases}
\]
To see this, any common divisor divides $2s^2$ and $2t^2$, and
$\gcd(s^2,t^2)=1$ therefore makes it divide $2$.  If the inputs
are both odd, all three coordinates are even; if their parities
differ, the first and third are odd.  They cannot both be even
because they are coprime.  Since
$|s^2-t^2|\leq s^2+t^2$ and $|2st|\leq s^2+t^2$, we obtain the
exact height formula
\begin{equation}\label{eq:circle-exact-height}
 \Ht(\Phi([s:t]))=\frac{s^2+t^2}{\epsilon}.
\end{equation}
For example $[2:1]$ yields $[3:4:5]$, while $[1:1]$ yields
$[0:1:1]$ after dividing the displayed coordinates by $2$.
The same identity is the elementary connection between rational
points on a conic and primitive Pythagorean triples.

\paragraph{An empty conic detected by a prime-power obstruction.}
The nonsingular conic
\[
 X^2+Y^2=3Z^2
\]
has no rational projective point.  Suppose it had one and choose
a primitive integer representative $(x,y,z)$.  The only square
residues modulo $3$ are $0$ and $1$, so $x^2+y^2\equiv0\pmod3$
forces $3\mid x$ and $3\mid y$.  Write $x=3u,y=3v$ and cancel
$3$ from the equation to obtain $z^2=3(u^2+v^2)$.  The same
square-residue table gives $3\mid z$, contradicting primitivity.
Every step is a finite divisibility calculation.

There \emph{are} projective points of this equation over the
field with three elements, for example $(0:0:1)$.  The obstruction
is the failure of such a residue class to lift to a solution modulo
$9$ with at least one coordinate not divisible by $3$.  Equivalently,
it is the two successive modulo-$3$ deductions just given.  This
distinction matters when formulating a correct finite local
certificate.

\subsection{Why these arguments fit the weak constructive theory}

No existence claim above asks for an unbounded search.  A seed point
is part of the hypotheses of the parametrization theorem; the
theorem does not search indefinitely for one.  A conic with no
supplied seed still has the bounded point enumerator, and a
specific local obstruction may certify emptiness as in the second
example.  A general local-to-global theorem for conics is a further
theorem, not an unstated step of this chapter.

The algebraic parts use fixed polynomial identities, gcd and
B\'ezout algorithms, and decidable branches.  The counting parts
only use bounded ranges.  Their induction invariants are equalities
and inequalities of primitive recursive counting functions; a
bounded finite-union assertion is represented by its primitive
recursive characteristic function.  Finite telescoping is proved
by quantifier-free induction, after clearing the positive
denominators or maintaining a rational pair code.  Distinctness of
the produced list is also a bounded predicate.

Thus the internal assertion of constructive infinitude is the
explicit arithmetic statement that for every input $n$, a specified
primitive recursive list has length $n+1$, distinct valid entries,
and the displayed height bound.  It does not quantify over an
infinite set, invoke a choice principle, or require induction on
an arbitrary formula.  The bijection and the counting estimates
are uniform first-order arithmetic theorems about the supplied
coefficient and point codes.
